\cvsection{Expériences professionnelles}

\begin{cventries}
  \cventry
    {Stagiaire ingénieur de recherche en IA}
    {INRIA, équipe Mnémosyne}
    {Bordeaux, France}
    {Mars 2026 -- Septembre 2026}
    {
      \begin{cvitems}
        \item {Conception et développement en autonomie, avec le suivi de mes encadrants, d’une application web répondant au besoin d’aider les chercheurs à trouver des articles scientifiques.}
        \item {Développement du backend en Python/FastAPI, d’API REST documentées avec OpenAPI et de la persistance MongoDB.}
        \item {Développement d’une interface React/Next.js permettant aux chercheurs d’utiliser l’application et ses différents outils.}
        \item {Conteneurisation de composants avec Docker et mise en place d’un pipeline CI/CD sur GitHub.}
        \item {Création d’un outil Python de collecte automatisée de données et de webhooks transmettant les résultats de tâches planifiées.}
        \item {Prise en charge d’une application web conteneurisée utilisant des modèles installés localement sur un Mac Studio, incluant leur installation et leur configuration.}
        \item {Évaluation d’OpenClaw dans un environnement Docker isolé avec une attention portée aux risques liés à l’accès autonome aux outils et aux fichiers.}
        \item {Définition d’un protocole expérimental, analyse de résultats et contribution à un article scientifique comparant plusieurs modèles.}
      \end{cvitems}
    }

  \cventry
    {Stagiaire de recherche en apprentissage par renforcement}
    {Université Karunya}
    {Coimbatore, Tamil Nadu, Inde}
    {Juin 2025 -- Sept. 2025}
    {
      \begin{cvitems}
        \item {Conception et comparaison de deux approches de navigation autonome pour un robot hospitalier.}
        \item {Déploiement et optimisation du logiciel et du modèle d’apprentissage par renforcement sur un Raspberry Pi embarqué.}
      \end{cvitems}
    }
\end{cventries}
